\subsection{Pre-registration of the selective-editing study}
\label{sec:supp_prereg}
The study was registered before the run in an internal document with dated revisions. \textbf{Criteria.} B1: fewer broken clips, false skips $\leq$2\% of the skipped clips, paired readability not lower than B0 by more than 0.5 points. B2: collateral change decreases ($p<0.05$), readability not lower than B0 by more than one point, and the human artifact question favours B2. B3: infeasible flags confirmed by manual inspection; results reported with infeasible clips counted as failures (main) and excluded (reference). Components are attributed separately before B4 is considered. \textbf{Revisions during development} (on the development clips only, before the run) concerned the B2 glow model and the B3 boundary cues. \textbf{Revisions during the run} (before any result was seen, recorded at the time): B3 feasibility was changed to use the character density of the source, measured on the best reading over all frames; B1 received its fifth condition after two false skips were found while checking decisions (a correct word above garbled small text; a near-miss reading split across boxes). All 526 clips were then re-run with the final code. \textbf{Round two} (defined after the results, before any human data): a colour-completion variant of B2 and a whole-word-subset variant of B3 were tested on nine and six diagnosis clips, respectively, and were not run at scale. \textbf{Excluding infeasible clips} (reference): readability B1 92.0, B3 91.8; after manual review, replacing the non-confirmed B3 flags by B0 outputs gives 90.3\%.

\subsection{Pre-registration of the appearance modules}
\label{sec:supp_prereg_neg}
Both modules of \cref{sec:exp_neg} were registered before the full run in the same internal document format. \textbf{R1 criteria} (on the clips where R1 applies, outside its development clips): at most one clip less readable than B1; paired appearance difference with 95\% CI upper bound below zero; in a blind study, R1 preferred in more than half of the non-tied judgements with no more artifacts; coverage reported. Only all three together count as success. \textbf{B1f criteria} (all clips and the clips outside the development sets): paired readability difference not below $-1$ point; paired appearance difference with CI upper bound below zero; naturalness not worse in a later human study; meeting the appearance criterion alone is not sufficient. \textbf{Revisions during development}, on the three R1 development clips only: multi-frame consensus reading with per-letter stability, word-initial and word-final insertion and deletion, letter-by-letter layout at the source spacing instead of horizontal compression, a font-consistency gate on the kept letters, SAT fitted on the kept letters, and baseline-aligned segmentation. \textbf{Contamination.} The first run on the registered 30-clip evaluation sample applied R1 to 6 clips; these were then viewed while debugging, which revealed that LaMa continued the kept letters into the erased span (fixed by taking the middle background from the full-erase background), that appended text could leave its carrier (fixed with the carrier-extent estimate of B3), and that compressed letters merged into blobs (fixed by rendering letters individually). Because they had been inspected, the six clips were moved to the development set, and the evaluation of R1 was replaced by the full 526-clip pass with the final code; a first full pass started with the earlier code was cancelled and its outputs were discarded. \textbf{Outcome.} B1f fails the appearance criterion; R1 cannot be evaluated because it applies to two clips outside its development set. The choice of a glyph-agnostic appearance distance as the criterion for font and orientation changes was a design error; we report the registered criterion unchanged.

\subsection{Pre-registration of the material generator}
\label{sec:supp_prereg_matgen}
Registered before the pilot. \textbf{Gates.} R0 (data): at least 90\% of 40 pilot pairs render, and in 10 inspected pairs the two renders differ only in the text. R1 (synthetic counterfactual transfer, on held-out pairs whose strings do not occur in training): the model closes at least 0.5 of the gap between a flat paste and the rendered truth (the 1.3B refiner of \cref{sec:exp_neg}: $-0.06$); glyph IoU $\geq 0.85$; reported per material, with metal, neon and glass not negative. R2 (real clips, zero-shot): the benchmark clips that both systems repair, outside the development sets, with B1's erased background, B1's glyph alpha and the original as the only inputs; one evaluation run together with \method and \steer; appearance distance against B1 with CI upper bound below zero and mean at most $-0.30$, readability not more than one point below B1, BG PSNR $\geq 41$\,dB. R3 (all 526 clips): the R2 criteria plus the temporal limits and the second reader. No benchmark clip is used for training or model selection. \textbf{Deviation.} R1 closed 0.482 against the registered 0.50 (IoU 0.954; every material positive); the gate rule called for a diagnosis before R2, and R2 was run anyway as a diagnostic, recorded as such before its result was seen. \textbf{Outcome.} R2 fails all three criteria (\cref{sec:exp_neg}); R3 was not run. The one capacity test registered after R1 (base width 96 to 128, 31M parameters, same data, seed and schedule) closes 0.498 on the full validation set, below the gate by 0.002, with the same ordering of materials (metal lowest at 0.34); it was never applied to a real clip. The test bounds this one registered step in width, not capacity in general. \textbf{Photometric ablation} (registered after the R2 result, before its run). The frozen weights are applied to the same 129 clips with (P1) one per-clip, per-channel affine map that takes the reference crop to the training-set mean and spread, applied to both inputs and inverted exactly on the predicted text layer; (P2) the predicted ink's per-channel mean and spread matched to the source text's, on non-emissive clips only (the front end's glow flag, or source ink brighter than its surround by more than 30 luminance levels), with gains clipped to $[0.5, 2]$; and both (P12). The apply step now composites the erased crop only inside the erase mask (the R2 run replaced the whole window, which is why its BG PSNR is 38.8\,dB for the model and the flat paste alike); the flat paste and the unaligned model (P0) are regenerated with the same code, and every row is evaluated in one run. Decision rules, in order: (a) a variant meets the R2 criteria against B1: photometry is the main cause; (b) its appearance distance is below P0's (CI upper bound below zero) with readability not worse: photometry is one cause and the remainder lies in the model; (c) no variant improves on P0: photometry is not the cause and the synthetic-to-real wording is withdrawn. Independently, (d) P2 against the flat paste asks whether the model's spatial structure adds anything beyond colour statistics. Retraining with photometric augmentation is deferred unless (a) or (b) holds. \textbf{Result} (\cref{tab:r2p}; one run, 129 clips, of which \rpNonem{} are non-emissive under the registered rule and \rpEmis{} are skipped by P2). P2 and P12 fall under (b): on the non-emissive clips the appearance difference to P0 is \rpPTwoNeMean{} [\rpPTwoNeLo, \rpPTwoNeHi], with \rpPTwoNeWins{} readings gained and \rpPTwoNeLosses{} lost ($p=\rpPTwoNeP$), and the difference to \method shrinks from \rpPZeroNeVsBMean{} [\rpPZeroNeVsBLo, \rpPZeroNeVsBHi] to \rpPTwoNeVsBMean{} [\rpPTwoNeVsBLo, \rpPTwoNeVsBHi]; (a) is not met, because readability stays below \method's (\rpPTwoNeRead\% vs.\ 100\%). P1 falls under (c) on all clips (\rpPOneAllMean{} [\rpPOneAllLo, \rpPOneAllHi]): it helps the non-emissive clips (\rpPOneNeMean{} [\rpPOneNeLo, \rpPOneNeHi]) and hurts the dark and emissive ones, as expected of a training set whose crops have a mean intensity of 213 of 255 and whose text differs from its surround by a median of 18 luminance levels (71\% of the training text by less than 30). Rule (d) gives a stronger answer than the registered one: P2 is not level with the flat paste but worse (\rpPTwoNeVsFlatMean{} [\rpPTwoNeVsFlatLo, \rpPTwoNeVsFlatHi]), and the flat paste, which keeps only the source's median colour, is as close to the original as \method on this glyph-agnostic distance (\rpFlatAllMad{} vs.\ \rpBoneAllMad) while reading worse (\rpFlatAllRead\%) and flickering more (its per-frame warp is not \method's compositing). We keep the negative result, attribute part of it to photometry, and do not retrain: the registered gate measures colour statistics, which a constant colour satisfies, so a better generator would first need a criterion that rewards structure and a human study, neither of which this paper has. With the erase-mask compositing, BG PSNR is \rpPZeroAllBG\,dB for every generated row.
\IfFileExists{tables/tab_r2p.tex}{\begin{table}[t]
\centering\small\setlength{\tabcolsep}{4pt}
\caption{Photometric domain-alignment ablation of the material generator (frozen v0 weights, the 129 R2 clips, one evaluation run; appearance distance normalised within this run and not comparable to other tables). P1 maps a clip's reference crop to the training statistics before the model and undoes the map on its output; P2 matches the predicted ink's colour statistics to the source text's on non-emissive clips; the flat paste is the geometry pipeline alone. The original is the garbled input.}
\label{tab:r2p}
\begin{tabular}{lccccc}
\toprule
Row & Read\,$\uparrow$ & Appear.\,$\downarrow$ & BG\,$\uparrow$ & Flicker\,$\downarrow$ & Jitter\,$\downarrow$\\
\midrule
\multicolumn{6}{l}{\textit{All 129 R2 clips}}\\
Original & 0.0 & -- & -- & 2.07 & 9.00\\
\steer & 100.0 & 1.21 & 36.1 & 1.87 & 8.66\\
\method (B1) & 100.0 & 1.65 & 42.2 & 1.75 & 7.32\\
Flat paste (geometry only) & 89.1 & 1.58 & 41.4 & 2.70 & 9.94\\
Generator v0 (P0, as registered) & 86.8 & 2.42 & 41.3 & 2.63 & 9.45\\
+ input alignment (P1) & 87.6 & 2.52 & 41.3 & 2.51 & 9.09\\
+ output calibration (P2) & 91.5 & 2.09 & 41.3 & 2.65 & 9.57\\
+ both (P12) & 88.4 & 2.45 & 41.3 & 2.53 & 9.18\\
\midrule
\multicolumn{6}{l}{\textit{Non-emissive clips where P2 acts (50)}}\\
Original & 0.0 & -- & -- & 1.50 & 6.92\\
\steer & 100.0 & 1.67 & 36.1 & 1.26 & 6.52\\
\method (B1) & 100.0 & 1.82 & 42.3 & 1.25 & 5.92\\
Flat paste (geometry only) & 88.0 & 1.32 & 41.5 & 2.53 & 9.41\\
Generator v0 (P0, as registered) & 76.0 & 2.97 & 41.5 & 2.41 & 8.56\\
+ input alignment (P1) & 82.0 & 2.11 & 41.5 & 2.45 & 8.81\\
+ output calibration (P2) & 88.0 & 2.11 & 41.5 & 2.47 & 8.89\\
+ both (P12) & 84.0 & 2.00 & 41.5 & 2.49 & 9.04\\
\bottomrule
\end{tabular}
\end{table}
}{}

\subsection{Pre-registration of the verified selection}
\label{sec:supp_prereg_hybrid}
Registered before the run. \textbf{Candidates}: the frozen \method (B1), \steer seed 1 and seed 2 (the seed-2 outputs existed already). \textbf{Verifier}: InternVL2.5-8B with the evaluation protocol (same crop, 5 frames, at least 4 exact); it was chosen over the pipeline's own OCR because EasyOCR accepts only 193 of the 442 \steer outputs the first reader judges correct, which caps a hybrid at 92.8\%. \textbf{Decision order}: preservation gate $\to$ original; verifier reads seed 1 $\to$ its local composite; else seed 2 $\to$ its composite; else B1. No threshold is tuned. \textbf{Local compositing}: weight 1 inside the per-frame text box and up to 0.3 box heights outside, linear fade over 0.12 box heights, per-channel affine colour match on the annulus 0.42--0.92 box heights outside; the three parameters were set by eye on three development clips before the run. \textbf{Scoring}: the first reader (Qwen2.5-VL-7B) and, on English, Pixtral-12B; the verifier never scores. \textbf{Criteria}: readability at least 2 points above B1 with McNemar $p<0.05$; no more broken clips than B1; BG PSNR $\geq 41$\,dB; appearance distance at least 0.30 below B1; flicker $\leq 1.85$ and jitter $\leq 8.40$ (no worse than the worse parent); Pixtral readability not more than 1 point below B1's. \textbf{Outcome}: five criteria met; the temporal criterion missed (flicker \hVFlicker, jitter \hVJitter). \textbf{Post-hoc variant} (defined after seeing this): the colour match is estimated per frame but applied as one per-clip median gain and offset, so that it cannot vary over time; it was composited and evaluated in a second, separate run together with every other row, which is the run of \cref{tab:hybrid}; in that run the variant passes the flicker limit (\hVCFlicker) and misses jitter (\hVCJitter{} vs.\ 8.40), so the criterion remains unmet for both variants. The first run (pre-registered variants only) gave \hVRead\% readability and the same five-of-six verdict. \textbf{Diagnosis.} The composites are written by the same video writer as every \method output, whereas \steer's files are its own H.264 encodes. Re-encoding \steer's outputs with our writer, with no other change, raises their flicker by $+0.150$ [$+0.141$, $+0.159$] and their jitter by $+0.436$ [$+0.402$, $+0.468$] (paired over the 526 clips); the per-clip-matched composite is then indistinguishable from the re-encoded \steer output ($+0.000$ [$-0.002$, $+0.002$] flicker, $+0.004$ [$-0.005$, $+0.014$] jitter), and the per-frame-matched composite adds $+0.035$ and $+0.083$. The temporal criterion was set against \steer's own encode, so it compares two encoders rather than two systems; we leave the verdict as registered and report the cause.

